\scanpage{204}
ALGÈBRES de LIE SEMI-SIMPLES \mspage{(3)}

Tribu n° 59

et $\mathrm{q}:\mathrm{S}\longrightarrow\mathrm{W}$, et si on suppose que $\mathrm{q}(\mathrm{s})\mathrm{B}(\mathrm{v},\mathrm{v}')=\mathrm{B}(\mathrm{r}(\mathrm{s})\mathrm{v},\mathrm{v}')+\mathrm{B}(\mathrm{v},\mathrm{r}'(\mathrm{s})\mathrm{v}')$, (resp. $=\mathrm{B}(\mathrm{r}(\mathrm{s})\mathrm{v},\mathrm{r}'(\mathrm{s})\mathrm{v}')$, pour tout $\mathrm{s}$ de $\mathrm{S}$, alors $\mathrm{B}$ applique $\mathrm{V}_{\lambda}\times\mathrm{V}'_{\mu}$ dans $\mathrm{W}_{\lambda+\mu}$, (resp. dans $\mathrm{W}_{\lambda\mu}$) ; idem pour les indices supérieurs.

Ce n° absorde un certain nombre de choses du texte.

\uline{n° 1} - Vider p. 1, les l.9-7 bas, (c'est faux). La prop. 1, le lemme, la prop. 2, a) et la prop. 3 (a,c) vont au dépotoir, (noter que la prop. 3, c) absorbe le lemme. Ici le b) de la prop. 3 sera un cor. de la prop. 2, (réduite à son b)).

Dans la prop. 2, b), un erratum à Alg. VII permettra de parler des valeurs propres de $\mathrm{x}_{\mathrm M}$ même si elles ne sont pas dans $\mathrm{k}$. Pas besoin de supposer $\underline{\mathrm h}$ nilpotente à la p. 3, l.4 ; raisonner directement en utilisant Alg. VII. Dans (i), ne pas faire intervenir le quotient simple, (y'en a pas en caract. $\mathrm{p}$). Pour (ii), utiliser le dépotoir : la sous-algèbre engendrée par $\mathrm{x}$ opère sur $\mathrm{M}$, et sur $\underline{\mathrm h}$ par représ. adjointe ; on a $\underline{\mathrm h}\times\mathrm{M}\longrightarrow\mathrm{M}$.

La ``descente'' du bas de la p. 3 sera faite en Alg. VII ; donc un rappel suffit ici ; prendre une extension galoisienne qqconque au lieu de $\bar{\mathrm{k}}$.

Dans la prop. 4, à la l.13 p. 5, dire : \ldots\ est $\underline{\mathrm{g}}$-invariante, (identité de Jacobi), donc $\underline{\mathrm{h}}$-invariante. Dire ``$\underline{\mathrm{h}}$-module trivial'' à la l.15. Dire ``appliqué sur $0$ par la forme ci-dessus'' à la l.16. Plus bas la sommation est incorrecte : sommer sur la moitié des $\alpha$. Vider la remarque du bas de la p. 5.

On vide la prop. 5 qui ne sert à rien. Une remarque donnera la nilpotence de $\operatorname{Ad}(\mathrm{x})$ sous les conditions de la prop. 4.

Remonter ici le résultat suivant, (venant du \S\ 5, n° 3) : soient $\mathrm{G}$ un groupe analytique, $\mathrm{n}$ le rang de son algèbre de Lie ; alors il existe un sous-groupe ouvert $\mathrm{U}$ de $\mathrm{G}$ tel que, pour tout $\mathrm{x}\in\mathrm{U}$, $\operatorname{Ad}(\mathrm{x})$ admette $1$ comme valeur propre de

% archive sheet number: 203
